\input{../header}
\usepackage{pgfplots}
\rhead{Your name: \rule{8cm}{0.15mm}}

\everymath{\displaystyle}
\begin{document}
%


%\onehalfspacing
\allowdisplaybreaks
%##################################################################
\section{Chapter 2 checkpoint!}

Chapter 1 scorecard:
\begin{center}
    \begin{tabular}{|m{3.75cm}|*{5}{m{2cm}|}} \hline
        Learning target: & DF1 & DF2 & DFa & DFb & AD2 \\\hline
        Your confidence level before starting (0-5): & &&&&\\\hline
        Your confidence level after the quiz (0-5): & &&&&\\\hline
        The mark you earned on this attempt: 
        & Success! \newline Try again!
        & Success! \newline Try again!
        & Success! \newline Try again!
        & Success! \newline Try again!
        & Success! \newline Try again! \\\hline

    \end{tabular}
\end{center}
Chapter 2 scorecard:
\begin{center}
    \begin{tabular}{|m{3.75cm}|*{5}{m{2cm}|}} \hline
        Learning target: & 2.1 & 2.3 & 2.5 & 2.7 & 2.all  \\\hline
        Your confidence level before starting (0-5): & &&&&\\\hline
        Your confidence level after the quiz (0-5): & &&&&\\\hline
        The mark you earned on this attempt: 
        & Success! \newline Try again!
        & Success! \newline Try again!
        & Success! \newline Try again!
        & Success! \newline Try again!
        & Success! \newline Try again! \\\hline

    \end{tabular}
\end{center}

Before anything else, please do the following:
\begin{itemize}
    \item Rank your confidence from 0-5 on each of the learning targets. 5 means ``I could teach a whole class about this;'' 0 means ``I am genuinely not sure I have heard these words before.''
    \item Write your name on this page and on each of the other pages of the quiz.
\end{itemize}

Then do the quiz! Some reminders:
\begin{itemize}
    \item Open notes, closed computer.
    \item If you need more room to write, use the back of the same learning target page, or ask me for some scratch paper.
    \item Read the questions carefully and make sure you're answering each part.
    \item Use good grammar for derivative problems.
    \item Show all your work and explain all your thinking!
\end{itemize}

When you are done:
\begin{itemize}
    \item Rank your confidence from 0-5 on each of the learning targets. 5 means ``I absolutely nailed that question for sure;'' 0 means ``oof, I definitely didn't get that one.''
    \item Make double sure your name is on every page, including any scratch paper.
    \item Hand in your work, separated by learning target.
\end{itemize}

Have fun and do your best! I believe in u $\heartsuit$


%%%%%%%%%%%%%%%%%%%%%%%%%%%%%%%%%%%%%%%%%%%%%%%%%%%%%%%%%
\pagebreak
%%%%%%%%%%%%%%%%%%%%%%%%%%%%%%%%%%%%%%%%%%%%%%%%%%%%%%%%%

\section{Learning target 2.1, version 1}

Compute derivatives using basic function rules, the constant multiple rule, and the sum rule.

\hrulefill

Demonstrate and explain how to find the derivative of the following functions. Be sure to write down which derivative rules (constant multiple, sum/difference, etc.) you are using in your work.

\begin{enumerate}
    \item \(f(x) = 2 \, \cos\left(x\right) + 4 \, e^{x}\) 
    \vfill
    \item \(h(t) = \sqrt[3]{t^4} + \frac{4}{t^{4}}\) \quad (Don't use the quotient rule or the chain rule.)
    \vfill
    \item \(g(w) = 6 \, w^{3} - 2 \, w^{2} - 9 \, w - 2\)
    \vfill
    \item \(C(t) = 64(e^7 + 7^e)- \dfrac{\sqrt[27]{193262}}{16-\pi^2}\cdot\sin(4)\) \quad
(Hint: Look carefully!)
\end{enumerate}

%%%%%%%%%%%%%%%%%%%%%%%%%%%%%%%%%%%%%%%%%%%%%%%%%%%%%%%%%
\pagebreak
%%%%%%%%%%%%%%%%%%%%%%%%%%%%%%%%%%%%%%%%%%%%%%%%%%%%%%%%%

\section{Learning target 2.3, version 1}

Compute derivatives using the product and quotient rules.

\hrulefill

\begin{enumerate}[leftmargin=15pt]
    \item Consider the function $g(t) = 7\cdot t^4$.
    \begin{enumerate}
        \item Compute $g'(t)$ using the product rule. Simplify as much as possible.
        \vfill
        \item Now compute $g'(t)$ again, \textit{without} the product rule.
        \vfill
        \item Are both ways \textit{correct}? Which way do you \textit{prefer}?
        \vfill
    \end{enumerate}
    \item Consider the function $f(t) = \dfrac{t^4}{7}$. 

    \begin{enumerate}
        \item Compute $f'(t)$ using the quotient rule. Simplify as much as possible.
        \vfill
        \item Now compute $f'(t)$ again, \textit{without} the quotient rule.\\
        (Hint: What's another way to say ``divide by 7''?)
        \vfill
        \item Are both ways \textit{correct}? Which way do you \textit{prefer}?
    \end{enumerate}
\end{enumerate}

%%%%%%%%%%%%%%%%%%%%%%%%%%%%%%%%%%%%%%%%%%%%%%%%%%%%%%%%%
\pagebreak
%%%%%%%%%%%%%%%%%%%%%%%%%%%%%%%%%%%%%%%%%%%%%%%%%%%%%%%%%

\section{Learning target 2.5, version 1}

Compute derivatives using the chain rule.

\hrulefill

Demonstrate and explain how to find the derivative of the following functions. Be sure to write down which derivative rules (chain, product, quotient, sum and difference, etc.) you are using.

\begin{enumerate}
    \item \(h(t)= 16 \, {\left(t + e^{t} + 2\right)}^{4}\)
    \vfill
    \item \(f(w)= 5 \, \sin\left(w^{7/3}\right)\)
    \vfill
    \item \(g(w)= 5 \, \Big(\sin\left(w\right)\Big)^{7/3}\)
    \vfill
    \item \(p(y) = 2\ln\left(\sin(y) + y^2 + 3\right)\)
    \vfill
\end{enumerate}

%%%%%%%%%%%%%%%%%%%%%%%%%%%%%%%%%%%%%%%%%%%%%%%%%%%%%%%%%
\pagebreak
%%%%%%%%%%%%%%%%%%%%%%%%%%%%%%%%%%%%%%%%%%%%%%%%%%%%%%%%%

\section{Learning target 2.all, version 1}

Compute derivatives of complicated functions using a combination of derivative rules.

\hrulefill

Demonstrate and explain how to find the derivative of the following functions. Be sure to write down which derivative rules (chain, product, quotient, sum and difference, etc.) you are using.

\begin{enumerate}
    \item \(g(x) = \sqrt{\cos\left(-2 \, x^{4} + 5\right)}\)
    \vfill
    \item \(f(w) = \left( \frac{6 \, w^{6} - 1}{5 \, w^2 + 3} \right)^{ 4 }\)
    \vfill
    \item \(h(y) = {\left(3 \, y^{2} + 7 \, y\right)}^{4}\cdot y^{1/4}\)
    \vfill
\end{enumerate}

%%%%%%%%%%%%%%%%%%%%%%%%%%%%%%%%%%%%%%%%%%%%%%%%%%%%%%%%%
\pagebreak
%%%%%%%%%%%%%%%%%%%%%%%%%%%%%%%%%%%%%%%%%%%%%%%%%%%%%%%%%

\section{Learning target 2.7, version 1}

Compute derivatives of implicitly-defined functions.

\hrulefill

Consider the curve defined implicitly by the equation
\[y^3-x\cdot y = -6.\]
	
\begin{enumerate}[leftmargin=15pt]

	\item Use implicit differentiation to find an expression for $\frac{dy}{dx}$ that depends on $x$ and $y$.  Show all of your work and thinking clearly, using correct notation.

	\vfill

	\item The point $(7,2)$ lies on this curve. (I promise this is true, but how would you check?) \\
	Find the exact slope of the tangent line to this curve at the point \((7,2)\). Show clearly how you determined this value, and label your work with proper notation.

	\vspace{1.25in}

	
	\item (Also LT 1.8!) Figure out the equation of the tangent line $L(x)$ to this curve at the point $(7,2)$, and use your tangent line equation to make a reasonable guess about the $y$-coordinate of the point $(6.8, \rule{0.75cm}{0.15mm} )$.
	\vfill

\end{enumerate}

\end{document}